\documentclass[../../main]{subfiles}
\begin{document}

This thesis investigated the feasibility of employing LLMs as decision components for solver selection in CP. The central research question concerned whether a general-purpose LLM, without task-specific fine-tuning, can infer appropriate solver choices from structural characteristics of combinatorial models within a predefined solver portfolio.

The study focused on MiniZinc~\cite{MiniZinc} models derived from benchmark instances of the MiniZinc Challenge~\cite{MznProblems}. Several input representations were evaluated, including raw MiniZinc scripts, structured feature vectors, natural-language descriptions, and natural-language abstractions derived from FlatZinc models. These representations were combined with different prompting strategies and inference configurations in order to assess their impact on solver selection performance. The objective was to determine whether structural information about a problem instance can be effectively interpreted by an LLM and translated into meaningful solver recommendations.

The experimental results indicate that general-purpose LLMs exhibit a limited but non-negligible capacity to interpret structural properties of constraint models. Initial experiments showed that na\"ive prompting strategies based on raw MiniZinc scripts produced results below the performance of the single best solver. Introducing structured natural-language descriptions produced modest improvements, suggesting that the representation of the input problem significantly influences the model's ability to interpret relevant structural information.

Further experiments explored more structured input formats. Feature vectors generated by the \texttt{mzn2feat} tool~\cite{mzn2feat} were evaluated both independently and in combination with other representations. However, this approach did not yield improvements over the best configurations obtained with natural-language inputs. Additional tuning of inference parameters, particularly sampling temperature, produced small performance gains, indicating that controlled stochasticity can influence solver selection behaviour.

Despite these adjustments, the overall performance of the LLM-based approach remained below that of the best single solver. A recurring pattern emerged in which the model consistently selected a small subset of globally dominant solvers across most instances. While this strategy is not inherently ineffective given the strong empirical performance of those solvers, it limits the discriminative capacity of the selection process and reduces sensitivity to instance-specific characteristics.

To further investigate this behaviour, additional experiments were conducted using a restricted solver portfolio. In this setting, accurate discrimination among solvers becomes necessary to achieve competitive performance. Under these conditions, the limitations of the approach became more apparent, as the LLM-based selection consistently performed below the best solver within the restricted set.

Motivated by the observed influence of natural-language representations, this thesis introduced \texttt{fzn2nl}~\cite{fzn2nl}, a tool that converts FlatZinc models into deterministic natural-language descriptions. The resulting representation preserves structural information while reducing syntactic noise and token usage. Experimental evaluation shows that this representation achieves performance comparable to other input modalities while requiring substantially fewer tokens, demonstrating improved efficiency in terms of context usage.

Additional experiments examined the internal decision dynamics of the LLM. Controlled perturbation tests revealed that solver names strongly influence the model's choices: when solver descriptions were permuted while preserving solver names, the resulting selections remained largely unchanged. Conversely, anonymizing solver names increased the influence of descriptive content, indicating that the model can partially associate solver characteristics with problem structure when name-based biases are removed. Nevertheless, solver choices remained concentrated on a small subset of candidates.

Qualitative analysis of generated explanations further indicates that the model possesses a coherent high-level understanding of constraint models and solver paradigms. The LLM correctly identifies structural aspects such as constraint types, objective structure, and variable domains, and it is generally capable of associating these properties with broad solver categories (e.g., CP, MILP, or CP-SAT). However, the model shows limited ability to discriminate between solvers within the same category or to relate fine-grained instance characteristics to empirically optimal solver choices.

The conclusions of this study must be interpreted in light of several limitations. The experimental campaign was constrained by API rate limits and token throughput restrictions imposed by commercial LLM services. These constraints limited the number of evaluated instances and prevented extensive repeated trials across different inference configurations.

A further limitation concerns the intrinsic non-determinism of LLM inference. Since model outputs depend on stochastic sampling mechanisms, identical prompts may produce different solver recommendations across runs. This variability affects the stability of individual predictions and complicates systematic evaluation.

In addition, the empirical analysis was restricted to benchmark instances derived from the MiniZinc Challenge. Although this benchmark suite represents a widely used evaluation standard in CP, it does not cover the full diversity of real-world modeling scenarios.

Future work may address these limitations by extending the evaluation to larger and more heterogeneous benchmark collections and by assessing more advanced LLM architectures with larger context windows and improved reasoning capabilities. Another promising direction concerns the integration of task-specific learning. Fine-tuning LLMs on solver-performance datasets could enable models to internalize empirical solver competitiveness and strengthen the relationship between structural instance features and solver selection decisions.

Overall, the results indicate that current general-purpose LLMs are not yet competitive with specialized solver-selection methods. Nevertheless, they demonstrate a meaningful capacity to interpret structural properties of constraint models and solver paradigms. These findings suggest that, with appropriate methodological developments and domain-specific adaptation, LLM-based components may become a viable element of future CP workflows.

\end{document}